Los requisitos del capítulo 02 se reproducen en la Tabla~\ref{tab:rt-02}.\par
\begin{tablalafrox}{Requisitos transversales, capítulo 02. Fuente: Bases Técnicas Transversales.}{tab:rt-02}%
  {L{1.8cm} Y L{2cm} L{1.4cm} L{2.2cm} L{1.5cm} L{2cm}}%
  {\cab{Código} & \cab{Exigencia} & \cab{Carácter} & \cab{Cumple} & \cab{Componente} & \cab{Sección} & \cab{Evidencia}}
  RT-02.01 & La solución se organizará en las ocho capas del numeral 2.1. El PROPONENTE presentará el diagrama de la arquitectura lógica identificando cada capa, sus componentes y las interfaces entre ellas. & Obligatorio &  &  &  &  \\
  RT-02.02 & La arquitectura será modular, con límites de contexto explícitos y acoplamiento débil. Se rechazará toda arquitectura monolítica que no permita desplegar de forma independiente sus componentes críticos. & Obligatorio &  &  &  &  \\
  RT-02.03 & La descripción de la arquitectura se ajustará a ISO/IEC/IEEE 42010, con vistas lógica, de procesos, de despliegue, de datos y de seguridad. & Obligatorio &  &  &  &  \\
  RT-02.04 & El PROPONENTE mantendrá un registro de decisiones de arquitectura (ADR) fechado, con la alternativa escogida, las alternativas descartadas y el criterio de decisión. El registro es entregable contractual y se actualizará durante toda la ejecución. & Obligatorio &  &  &  &  \\
  RT-02.05 & La capa de servicios de negocio será sin estado. El estado de sesión y el estado de proceso residirán en almacenes externos con alta disponibilidad. & Obligatorio &  &  &  &  \\
  RT-02.06 & Toda operación de escritura expuesta a reintentos será idempotente, con clave de idempotencia declarada por el cliente y ventana de deduplicación documentada. & Obligatorio &  &  &  &  \\
  RT-02.07 & Los flujos de eventos garantizarán entrega al menos una vez, con deduplicación en el consumidor y orden garantizado dentro de la partición o del agregado cuando el proceso lo exija. & Obligatorio &  &  &  &  \\
  RT-02.08 & La solución implementará patrones de resiliencia demostrables: reintento con retroceso exponencial y variación aleatoria, cortacircuitos, mamparos de aislamiento, límites de tasa y tiempo de espera explícito en toda llamada remota. No se admiten llamadas remotas sin tiempo de espera. & Obligatorio &  &  &  &  \\
  RT-02.09 & La solución degradará de forma elegante: ante la indisponibilidad de un componente no crítico deberá continuar operando en modo reducido, informando la degradación a la persona usuaria, y nunca fallar de forma total. & Obligatorio &  &  &  &  \\
  RT-02.10 & Las capas de aplicación e integración escalarán horizontalmente de forma automática, con umbrales, límites superiores y costo asociado declarados en la oferta. & Obligatorio &  &  &  &  \\
  RT-02.11 & El PROPONENTE declarará explícitamente los puntos únicos de falla que subsistan en su arquitectura y justificará por qué son aceptables. Omitir esta declaración cuando existan puntos únicos de falla se evaluará como observación grave. & Obligatorio &  &  &  &  \\
  RT-02.12 & La solución admitirá su replicación a nuevas unidades, sitios, sucursales o filiales del CLIENTE sin rediseño arquitectónico, mediante parametrización o multi-tenencia. & Según caso &  &  &  &  \\
  RT-02.13 & El PROPONENTE presentará un modelo de dominio del negocio del caso, con las entidades principales, sus relaciones y los eventos de negocio que las modifican. & Obligatorio &  &  &  &  \\
  RT-02.14 & Se valorará la aplicación documentada de patrones de arquitectura evolutiva que permitan sustituir un componente sin reescribir la solución: capa anticorrupción frente a sistemas heredados, estrangulamiento progresivo y abstracción de proveedores. & Deseable &  &  &  &  \\
\end{tablalafrox}
Los requisitos del capítulo 03 se reproducen en la Tabla~\ref{tab:rt-03}.\par
\begin{tablalafrox}{Requisitos transversales, capítulo 03. Fuente: Bases Técnicas Transversales.}{tab:rt-03}%
  {L{1.8cm} Y L{2cm} L{1.4cm} L{2.2cm} L{1.5cm} L{2cm}}%
  {\cab{Código} & \cab{Exigencia} & \cab{Carácter} & \cab{Cumple} & \cab{Componente} & \cab{Sección} & \cab{Evidencia}}
  RT-03.01 & El PROPONENTE declarará el proveedor de nube pública, la región primaria y la región secundaria utilizadas. El proveedor deberá contar con presencia de región o zona en Chile o en Sudamérica. & Obligatorio &  &  &  &  \\
  RT-03.02 & Todos los componentes con requisito de alta disponibilidad se desplegarán en al menos dos zonas de disponibilidad. No se aceptará un diseño en una sola zona. & Obligatorio &  &  &  &  \\
  RT-03.03 & La totalidad de la infraestructura se definirá como código, versionada en el repositorio del CLIENTE, revisable y reproducible. No se admite infraestructura creada manualmente por consola, salvo la cuenta raíz inicial, cuya creación deberá documentarse. & Obligatorio &  &  &  &  \\
  RT-03.04 & La red se segmentará por capas, con subredes privadas para aplicación y datos, y exposición pública restringida a la capa de borde. Ningún componente de datos será alcanzable desde Internet. & Obligatorio &  &  &  &  \\
  RT-03.05 & El PROPONENTE privilegiará servicios administrados por sobre servicios autoadministrados cuando ello reduzca el riesgo operacional, y justificará cada excepción. & Obligatorio &  &  &  &  \\
  RT-03.06 & Se aplicarán prácticas FinOps: etiquetado obligatorio de todos los recursos por ambiente, módulo y centro de costo; presupuestos con alertas de desviación; y reporte mensual de consumo desglosado entregado al CLIENTE. & Obligatorio &  &  &  &  \\
  RT-03.07 & El PROPONENTE declarará su estrategia de reversibilidad y de mitigación del bloqueo por proveedor, identificando qué componentes son portables, cuáles no lo son y cuál sería el esfuerzo estimado de una migración. & Obligatorio &  &  &  &  \\
  RT-03.08 & La solución empleará instancias reservadas, planes de ahorro o capacidad comprometida cuando el perfil de carga lo justifique, y lo reflejará en la estructura de costos. & Deseable &  &  &  &  \\
  RT-03.09 & Se valorará el uso de cómputo sin servidor o de contenedores administrados para las cargas de perfil variable, con el análisis comparativo de costo frente a instancias permanentes. & Deseable &  &  &  &  \\
  RT-03.10 & El componente on-premise operará de forma autónoma y degradada ante la pérdida total del enlace con la nube, durante un período mínimo de 24 horas continuas o el mayor que fije el caso. & Obligatorio &  &  &  &  \\
  RT-03.11 & Durante la operación desconectada, la solución continuará registrando las transacciones operacionales críticas de forma local, con integridad garantizada y sin pérdida de datos. & Obligatorio &  &  &  &  \\
  RT-03.12 & Restablecido el enlace, la sincronización será automática, con reconciliación determinista de conflictos, regla de resolución documentada y bitácora auditable de las decisiones aplicadas. & Obligatorio &  &  &  &  \\
  RT-03.13 & El PROPONENTE declarará qué funciones NO estarán disponibles en modo desconectado y qué procedimiento manual las suple. La ausencia de esta declaración se evaluará como observación grave. & Obligatorio &  &  &  &  \\
  RT-03.14 & Los equipos on-premise críticos serán redundantes. El almacenamiento local tolerará la falla de al menos un disco; el PROPONENTE declarará el nivel RAID escogido y lo justificará frente a las alternativas. & Obligatorio &  &  &  &  \\
  RT-03.15 & Los sistemas on-premise se endurecerán conforme a los CIS Benchmarks aplicables, con gestión centralizada de parches y ventana de aplicación acordada con el CLIENTE. & Obligatorio &  &  &  &  \\
  RT-03.16 & El monitoreo del componente on-premise se integrará a la misma plataforma de observabilidad que la nube, con alertamiento unificado. & Obligatorio &  &  &  &  \\
  RT-03.17 & El enlace entre el sitio on-premise y la nube será redundante, con caminos físicos y proveedores distintos, y conmutación automática con tiempo de conmutación declarado. & Obligatorio &  &  &  &  \\
  RT-03.18 & Los dispositivos de borde y de terreno se administrarán de forma remota y centralizada: inventario, configuración, actualización de firmware y de aplicación, bloqueo y borrado remoto. & Obligatorio &  &  &  &  \\
  RT-03.19 & Se valorará el procesamiento en el borde de las cargas que lo admitan ---filtrado, agregación previa, inferencia local--- reduciendo el volumen transferido y la dependencia del enlace. & Deseable &  &  &  &  \\
  RT-03.20 & El PROPONENTE dimensionará el ancho de banda requerido por sitio, en régimen normal y en peak, y lo justificará con el cálculo de volumen de transacciones y de datos. & Obligatorio &  &  &  &  \\
  RT-03.21 & La conexión entre la red del CLIENTE y la nube se establecerá mediante enlace privado dedicado o red privada virtual con cifrado, según lo que el volumen y la criticidad justifiquen. & Obligatorio &  &  &  &  \\
  RT-03.22 & El acceso remoto de las personas trabajadoras del CLIENTE, incluido el trabajo desde el hogar, se resolverá con acceso a la red de confianza cero, con verificación de postura del dispositivo. No se admite exponer servicios internos directamente a Internet. & Obligatorio &  &  &  &  \\
  RT-03.23 & La red inalámbrica de los sitios operacionales, cuando el caso la requiera, contará con segmentación por tipo de dispositivo, autenticación por certificado o credencial de empresa y cobertura verificada mediante estudio de sitio. & Según caso &  &  &  &  \\
  RT-03.24 & El PROPONENTE declarará la calidad de servicio y la priorización de tráfico aplicada a las transacciones operacionales críticas frente al tráfico administrativo. & Deseable &  &  &  &  \\
\end{tablalafrox}
Los requisitos del capítulo 04 se reproducen en la Tabla~\ref{tab:rt-04}.\par
\begin{tablalafrox}{Requisitos transversales, capítulo 04. Fuente: Bases Técnicas Transversales.}{tab:rt-04}%
  {L{1.8cm} Y L{2cm} L{1.4cm} L{2.2cm} L{1.5cm} L{2cm}}%
  {\cab{Código} & \cab{Exigencia} & \cab{Carácter} & \cab{Cumple} & \cab{Componente} & \cab{Sección} & \cab{Evidencia}}
  RT-04.01 & Los cinco ambientes del numeral 4.1 estarán habilitados y operativos como condición del hito H3 del Formulario E-25. & Obligatorio &  &  &  &  \\
  RT-04.02 & Preproducción será equivalente a producción en topología, versiones de componentes y configuración. Las diferencias que subsistan por costo se declararán expresamente y se justificarán. & Obligatorio &  &  &  &  \\
  RT-04.03 & El código residirá en un sistema de control de versiones con ramas protegidas, revisión obligatoria por pares y prohibición de escritura directa sobre la rama principal. & Obligatorio &  &  &  &  \\
  RT-04.04 & Existirá trazabilidad completa entre requerimiento, incidencia, cambio de código, prueba ejecutada y despliegue realizado. & Obligatorio &  &  &  &  \\
  RT-04.05 & El flujo de integración continua ejecutará, como mínimo: compilación, pruebas unitarias, análisis estático de código, análisis de composición de software, escaneo de secretos y escaneo de imágenes de contenedor, con criterios de bloqueo automático del despliegue. & Obligatorio &  &  &  &  \\
  RT-04.06 & Los despliegues serán automatizados y reproducibles, con reversión automatizada y sin intervención manual en el paso a producción. & Obligatorio &  &  &  &  \\
  RT-04.07 & La estrategia de despliegue permitirá liberar sin interrupción del servicio: azul-verde, canario o despliegue progresivo. Se declarará cuál se emplea y se demostrará en Preproducción antes de cada paso a producción. & Obligatorio &  &  &  &  \\
  RT-04.08 & Toda configuración estará externalizada del artefacto y gestionada por ambiente. Un mismo artefacto deberá poder promoverse de QA a Preproducción y a Producción sin recompilación. & Obligatorio &  &  &  &  \\
  RT-04.09 & Los secretos residirán en un gestor de secretos con rotación automática y auditoría de acceso. Queda prohibida toda credencial embebida en código, imágenes o archivos de configuración. & Obligatorio &  &  &  &  \\
  RT-04.10 & Las migraciones de esquema de base de datos serán versionadas, reversibles y ejecutadas de forma automatizada, con estrategia de compatibilidad que permita convivir dos versiones de la aplicación durante el despliegue. & Obligatorio &  &  &  &  \\
  RT-04.11 & La cobertura de pruebas automatizadas del código de lógica de negocio será de al menos 70 \%, con umbral bloqueante en el flujo de integración continua. & Obligatorio &  &  &  &  \\
  RT-04.12 & El PROPONENTE declarará su frecuencia de despliegue objetivo, su tiempo desde el compromiso de código hasta producción, su tasa de cambios fallidos y su tiempo de restauración, y los medirá durante la Operación. & Obligatorio &  &  &  &  \\
  RT-04.13 & Los ambientes no productivos se apagarán o reducirán fuera del horario de uso, con el ahorro reflejado en la estructura de costos. & Deseable &  &  &  &  \\
  RT-04.14 & Se valorará la existencia de ambientes efímeros por rama o por incidencia, creados y destruidos automáticamente. & Deseable &  &  &  &  \\
\end{tablalafrox}
Los requisitos del capítulo 05 se reproducen en la Tabla~\ref{tab:rt-05}.\par
\begin{tablalafrox}{Requisitos transversales, capítulo 05. Fuente: Bases Técnicas Transversales.}{tab:rt-05}%
  {L{1.8cm} Y L{2cm} L{1.4cm} L{2.2cm} L{1.5cm} L{2cm}}%
  {\cab{Código} & \cab{Exigencia} & \cab{Carácter} & \cab{Cumple} & \cab{Componente} & \cab{Sección} & \cab{Evidencia}}
  RT-05.01 & El PROPONENTE entregará el modelo de datos documentado y un diccionario de datos con el nombre, el tipo, el dominio de valores, la obligatoriedad, el propietario y la sensibilidad de cada atributo. & Obligatorio &  &  &  &  \\
  RT-05.02 & El PROPONENTE justificará la selección del paradigma y del motor de persistencia: relacional o no relacional, garantías transaccionales, y la posición escogida entre consistencia y disponibilidad conforme al teorema CAP, para cada dominio de datos. & Obligatorio &  &  &  &  \\
  RT-05.03 & Toda operación de negocio será trazable: la solución permitirá reconstruir quién, qué, cuándo, desde qué dispositivo y con qué valores anteriores y posteriores, para cualquier registro y en cualquier momento del período de retención. & Obligatorio &  &  &  &  \\
  RT-05.04 & La calidad de datos se gestionará conforme a ISO/IEC 25012, con validación en el punto de captura, indicadores de completitud, exactitud y consistencia, y tablero de calidad disponible para el CLIENTE. & Obligatorio &  &  &  &  \\
  RT-05.05 & El almacenamiento transaccional y el analítico estarán separados. Ninguna consulta analítica podrá degradar el desempeño de la operación. & Obligatorio &  &  &  &  \\
  RT-05.06 & La solución permitirá exportar la totalidad de la información del CLIENTE en formatos abiertos y documentados, en cualquier momento del Contrato, sin costo adicional y sin intervención del ADJUDICATARIO. & Obligatorio &  &  &  &  \\
  RT-05.07 & El PROPONENTE declarará la política de retención, archivado y eliminación por cada dominio de datos, coherente con la normativa aplicable al caso, e implementará un procedimiento verificable de eliminación segura. & Obligatorio &  &  &  &  \\
  RT-05.08 & Los datos personales se tratarán conforme al Artículo 85° de las Bases Administrativas, con seudonimización o cifrado a nivel de campo para las categorías sensibles que el caso identifique. & Obligatorio &  &  &  &  \\
  RT-05.09 & El PROPONENTE presentará una estrategia de gestión de datos maestros que evite la duplicación de entidades compartidas entre módulos y con sistemas externos. & Obligatorio &  &  &  &  \\
  RT-05.10 & Se valorará la implementación de un catálogo de datos con linaje automatizado, que permita rastrear el origen de cada indicador de negocio hasta su fuente. & Deseable &  &  &  &  \\
  RT-05.11 & El PROPONENTE presentará un plan de migración con alcance, origen, volumen, reglas de transformación, criterios de calidad, estrategia de ejecución y plan de reversión. & Obligatorio &  &  &  &  \\
  RT-05.12 & La migración incluirá una etapa de perfilado y saneamiento previo, con informe de los defectos detectados en los datos de origen y la decisión adoptada sobre cada uno. & Obligatorio &  &  &  &  \\
  RT-05.13 & Se ejecutarán al menos dos ensayos completos de migración sobre Preproducción antes de la migración definitiva, con medición del tiempo total y del resultado de la conciliación. & Obligatorio &  &  &  &  \\
  RT-05.14 & La conciliación posterior a la migración será cuantitativa y verificable: recuentos, sumas de control y muestreo dirigido. Toda diferencia deberá quedar explicada. & Obligatorio &  &  &  &  \\
  RT-05.15 & Los datos históricos que no se migren quedarán accesibles en un repositorio de consulta durante el período de retención que fije el caso. & Según caso &  &  &  &  \\
  RT-05.16 & Los servicios síncronos se documentarán en OpenAPI 3.1 y los flujos dirigidos por eventos en AsyncAPI 2.6 o superior. La documentación se generará desde el código y se mantendrá actualizada automáticamente. & Obligatorio &  &  &  &  \\
  RT-05.17 & Los contratos de interfaz se versionarán semánticamente, con compatibilidad hacia atrás y política de obsolescencia con preaviso mínimo de seis meses. & Obligatorio &  &  &  &  \\
  RT-05.18 & La autenticación entre sistemas empleará OAuth 2.1 con credenciales de cliente o autenticación mutua TLS. Queda prohibida la autenticación por clave estática en la ruta de la dirección web. & Obligatorio &  &  &  &  \\
  RT-05.19 & Toda integración registrará la transacción de entrada y de salida, con identificador de correlación común que permita seguir una operación de negocio a través de todos los sistemas involucrados. & Obligatorio &  &  &  &  \\
  RT-05.20 & Las integraciones con sistemas heredados o de terceros se aislarán mediante una capa anticorrupción, de modo que un cambio en el sistema externo no propague su modelo al núcleo de la solución. & Obligatorio &  &  &  &  \\
  RT-05.21 & El PROPONENTE declarará, por cada integración, el modo (síncrono o asíncrono), el volumen esperado, la ventana de disponibilidad del sistema contraparte y el comportamiento de la solución cuando ese sistema no responde. & Obligatorio &  &  &  &  \\
  RT-05.22 & La solución soportará la carga y descarga masiva de información en formatos abiertos, con validación previa, informe de errores por registro y procesamiento parcial. & Obligatorio &  &  &  &  \\
  RT-05.23 & Se emplearán los estándares sectoriales de intercambio que las Bases Técnicas del caso identifiquen. & Según caso &  &  &  &  \\
  RT-05.24 & Se valorará la publicación de un portal de servicios para desarrolladores con documentación navegable, ambiente de pruebas y credenciales de prueba autoservidas. & Deseable &  &  &  &  \\
  RT-05.25 & La solución proveerá una capa analítica con tableros operacionales y de gestión, construidos sobre los indicadores que las Bases Técnicas del caso definan. & Obligatorio &  &  &  &  \\
  RT-05.26 & Los tableros permitirán filtrar por período, unidad organizacional y dimensiones propias del caso, y profundizar desde el indicador agregado hasta la transacción de origen. & Obligatorio &  &  &  &  \\
  RT-05.27 & El CLIENTE podrá construir sus propios informes sin intervención del ADJUDICATARIO, mediante una herramienta de autoservicio con modelo semántico documentado. & Obligatorio &  &  &  &  \\
  RT-05.28 & Todo informe será exportable en formatos abiertos y programable para envío automático por calendario. & Obligatorio &  &  &  &  \\
  RT-05.29 & La latencia máxima entre la ocurrencia de una transacción y su disponibilidad en la capa analítica será la que fije el caso y, en su defecto, no superará las 4 horas. & Según caso &  &  &  &  \\
  RT-05.30 & Se valorará la incorporación de analítica predictiva pertinente al proceso del caso, con el modelo, sus variables, su métrica de desempeño y su plan de reentrenamiento documentados. & Deseable &  &  &  &  \\
\end{tablalafrox}
Los requisitos del capítulo 06 se reproducen en la Tabla~\ref{tab:rt-06}.\par
\begin{tablalafrox}{Requisitos transversales, capítulo 06. Fuente: Bases Técnicas Transversales.}{tab:rt-06}%
  {L{1.8cm} Y L{2cm} L{1.4cm} L{2.2cm} L{1.5cm} L{2cm}}%
  {\cab{Código} & \cab{Exigencia} & \cab{Carácter} & \cab{Cumple} & \cab{Componente} & \cab{Sección} & \cab{Evidencia}}
  RT-06.01 & El espacio asignado será de uso exclusivo de la solución y estará aislado de otras dependencias del CLIENTE, con acceso independiente. & Obligatorio &  &  &  &  \\
  RT-06.02 & Los muros no estructurales del recinto contarán con blindaje perimetral; el PROPONENTE especificará el material y la resistencia. & Obligatorio &  &  &  &  \\
  RT-06.03 & El PROPONENTE entregará el plano de distribución interna del recinto, con la separación de las zonas de generadores, baterías, climatización, servidores, comunicaciones, trabajo y respaldo. & Obligatorio &  &  &  &  \\
  RT-06.04 & El piso técnico, la canalización, el cableado estructurado y el etiquetado se ejecutarán conforme a norma, con documentación de la certificación de cada enlace. & Obligatorio &  &  &  &  \\
  RT-06.05 & Los racks de servidores serán independientes de los racks de equipos de comunicación. Se declarará la ocupación proyectada de cada rack y su margen de crecimiento. & Obligatorio &  &  &  &  \\
  RT-06.06 & La obra civil de separación de las instalaciones es de cargo del CLIENTE; su especificación técnica y su coordinación son de cargo del PROPONENTE. & Obligatorio &  &  &  &  \\
  RT-06.07 & El suministro eléctrico de los equipos será ininterrumpido, con sistema de alimentación ininterrumpida dimensionado para una autonomía mínima de 30 minutos a plena carga. & Obligatorio &  &  &  &  \\
  RT-06.08 & La capacidad de generación autónoma asegurará un rango mínimo de 24 horas continuas de operación, con estanque de combustible dimensionado y contrato de reabastecimiento declarado. & Obligatorio &  &  &  &  \\
  RT-06.09 & La instalación eléctrica del recinto será independiente de la del resto del edificio y cumplirá la normativa eléctrica chilena vigente, incluida la NCh Elec. 2777 sobre sistemas de puesta a tierra. & Obligatorio &  &  &  &  \\
  RT-06.10 & Se efectuará revisión y medición semestral de las instalaciones eléctricas del recinto, con informe entregable al CLIENTE. & Obligatorio &  &  &  &  \\
  RT-06.11 & El PROPONENTE declarará la carga eléctrica proyectada en kW, el factor de potencia y la eficiencia en el uso de la energía (PUE) estimada del recinto. & Obligatorio &  &  &  &  \\
  RT-06.12 & Se valorará la redundancia de alimentación en configuración 2N o N+1 con doble acometida y transferencia automática. & Deseable &  &  &  &  \\
  RT-06.13 & El recinto contará con climatización de precisión para operación continua, redundante en configuración N+1, con control de temperatura y de humedad relativa dentro de los rangos que recomienda el fabricante del equipamiento. & Obligatorio &  &  &  &  \\
  RT-06.14 & Se monitorearán en línea la temperatura, la humedad y la presencia de agua, con alertamiento integrado a la plataforma de observabilidad. & Obligatorio &  &  &  &  \\
  RT-06.15 & El PROPONENTE declarará la estrategia de contención de pasillo frío o caliente y su efecto en la eficiencia energética. & Deseable &  &  &  &  \\
  RT-06.16 & El recinto contará con detección temprana por aspiración de aire con tecnología láser, tipo AnaLASER o equivalente de prestaciones iguales o superiores. & Obligatorio &  &  &  &  \\
  RT-06.17 & La extinción será automática mediante agente limpio tipo FM-200 o equivalente, con aprobación UL e instalación conforme a norma NFPA. & Obligatorio &  &  &  &  \\
  RT-06.18 & Se proveerá un sistema secundario de extintores portátiles habilitados, con mantención y certificación vigentes. & Obligatorio &  &  &  &  \\
  RT-06.19 & El sistema de detección y extinción se integrará al monitoreo en línea y notificará al NOC y a la contraparte del CLIENTE. & Obligatorio &  &  &  &  \\
  RT-06.20 & El ingreso al recinto se controlará mediante seguridad física y control de acceso biométrico basado principalmente en biometría facial, con AFIS como respaldo. Se admite proponer sistemas de mayor seguridad. & Obligatorio &  &  &  &  \\
  RT-06.21 & Todo ingreso y egreso quedará registrado en una bitácora auditable, con identificación de la persona, fecha, hora y motivo, conservada por el período de retención declarado. & Obligatorio &  &  &  &  \\
  RT-06.22 & Entre el acceso principal y el término del pasillo de la zona de control se dispondrá un espacio para la atención de personas en proceso de enrolamiento. Se evaluará mejor la existencia de una estación de enrolamiento fuera de las instalaciones del recinto técnico. & Obligatorio &  &  &  &  \\
  RT-06.23 & Al término del pasillo se instalará un acceso que impida el paso de más de una persona a la vez, con nueva verificación de identidad previa al ingreso. & Obligatorio &  &  &  &  \\
  RT-06.24 & El recinto contará con videovigilancia y monitoreo IP, con imágenes en línea y disponibles para visualización de al menos los últimos 30 días. Las grabaciones anteriores se respaldarán en un medio secundario recuperable y auditable. & Obligatorio &  &  &  &  \\
  RT-06.25 & El PROPONENTE declarará el procedimiento de acceso de terceros ---fabricantes, mantenedores, auditores--- con acompañamiento obligatorio y registro. & Obligatorio &  &  &  &  \\
  RT-06.26 & Se habilitará un servicio de custodia de medios de respaldo para el sitio primario, en un medio físico transportable a otro lugar cuando el CLIENTE lo determine. Se admite proponer una solución más segura y eficiente, debidamente justificada. & Obligatorio &  &  &  &  \\
  RT-06.27 & El recinto de custodia cumplirá exigencias de luminosidad, humedad, ventilación y cualquier otro factor que pueda afectar la calidad y la disponibilidad de los medios. & Obligatorio &  &  &  &  \\
  RT-06.28 & Se llevará un inventario de medios con rotación, verificación periódica de legibilidad y registro de todo movimiento de entrada y de salida. & Obligatorio &  &  &  &  \\
  RT-06.29 & El PROPONENTE habilitará el espacio físico necesario para el personal encargado de la operación y administración de la plataforma, con estaciones de trabajo, telefonía, conexión a Internet y todo elemento que permita realizar la labor en condiciones adecuadas. & Obligatorio &  &  &  &  \\
  RT-06.30 & El espacio de operación estará separado de la sala de equipos y no requerirá el ingreso al recinto técnico para las labores habituales de operación. & Obligatorio &  &  &  &  \\
  RT-06.31 & Las instalaciones sanitarias, las zonas de seguridad ante emergencia y las áreas exteriores existentes en el edificio del CLIENTE podrán utilizarse y no deben implementarse nuevamente. El PROPONENTE declarará de cuáles hará uso. & Obligatorio &  &  &  &  \\
  RT-06.32 & El acceso a las redes de comunicaciones estará provisto a través de rutas físicas distintas, con ingreso al edificio por puntos separados. & Obligatorio &  &  &  &  \\
  RT-06.33 & El PROPONENTE proveerá toda la conectividad, la seguridad y las canalizaciones, cañerías o ductos requeridos para cumplir los niveles de servicio comprometidos. & Obligatorio &  &  &  &  \\
  RT-06.34 & Se privilegiará al PROPONENTE que ofrezca o provea especificaciones nuevas o mejores que las aquí establecidas, debidamente fundamentadas. & Deseable &  &  &  &  \\
\end{tablalafrox}
Los requisitos del capítulo 07 se reproducen en la Tabla~\ref{tab:rt-07}.\par
\begin{tablalafrox}{Requisitos transversales, capítulo 07. Fuente: Bases Técnicas Transversales.}{tab:rt-07}%
  {L{1.8cm} Y L{2cm} L{1.4cm} L{2.2cm} L{1.5cm} L{2cm}}%
  {\cab{Código} & \cab{Exigencia} & \cab{Carácter} & \cab{Cumple} & \cab{Componente} & \cab{Sección} & \cab{Evidencia}}
  RT-07.01 & El PROPONENTE declarará la modalidad escogida ---activo-activo o activo-pasivo--- y la justificará frente al costo, al RTO comprometido y a la complejidad operacional que introduce. & Obligatorio &  &  &  &  \\
  RT-07.02 & El sitio secundario estará emplazado a una distancia suficiente del principal para no verse afectado por el mismo evento de fuerza mayor. El PROPONENTE declarará la distancia y el análisis de amenazas comunes considerado. & Obligatorio &  &  &  &  \\
  RT-07.03 & La replicación de datos será continua, con medición y alertamiento del retraso de replicación. & Obligatorio &  &  &  &  \\
  RT-07.04 & El objetivo de tiempo de recuperación (RTO) no superará 4 horas y el objetivo de punto de recuperación (RPO) no superará 15 minutos para los servicios críticos, salvo exigencia superior del caso. & Obligatorio &  &  &  &  \\
  RT-07.05 & El procedimiento de conmutación estará documentado, automatizado en la mayor medida posible y ejecutable por el personal del CLIENTE tras la transferencia de conocimiento. & Obligatorio &  &  &  &  \\
  RT-07.06 & Existirá un procedimiento de retorno al sitio principal igualmente documentado y probado, con reconciliación de los datos generados durante la contingencia. & Obligatorio &  &  &  &  \\
  RT-07.07 & El plan de recuperación ante desastres se probará al menos dos veces al año mediante conmutación real, con informe de resultados, medición del RTO y del RPO efectivamente alcanzados y plan de corrección de las brechas detectadas. & Obligatorio &  &  &  &  \\
  RT-07.08 & Se valorará que la conmutación sea automática ante la detección de indisponibilidad, con criterio de disparo declarado y protección contra conmutación innecesaria. & Deseable &  &  &  &  \\
  RT-07.09 & La política de respaldo seguirá el esquema 3-2-1-0: tres copias, en dos medios distintos, una fuera de sitio, una inmutable o fuera de línea y cero errores de verificación de restauración. & Obligatorio &  &  &  &  \\
  RT-07.10 & Los respaldos estarán cifrados en reposo y en tránsito, con clave gestionada de forma independiente de la infraestructura respaldada. & Obligatorio &  &  &  &  \\
  RT-07.11 & Las copias inmutables estarán protegidas contra borrado y contra modificación durante su período de retención, incluso frente a credenciales administrativas comprometidas. & Obligatorio &  &  &  &  \\
  RT-07.12 & Se ejecutará y documentará una prueba de restauración al menos mensual, sobre una muestra representativa, con medición del tiempo efectivo de restauración. & Obligatorio &  &  &  &  \\
  RT-07.13 & El PROPONENTE declarará, por cada dominio de datos, la frecuencia de respaldo, el período de retención y el tiempo estimado de restauración completa. & Obligatorio &  &  &  &  \\
  RT-07.14 & Los respaldos permitirán la restauración granular: un registro, una tabla, un módulo o el sistema completo. & Deseable &  &  &  &  \\
\end{tablalafrox}
Los requisitos del capítulo 08 se reproducen en la Tabla~\ref{tab:rt-08}.\par
\begin{tablalafrox}{Requisitos transversales, capítulo 08. Fuente: Bases Técnicas Transversales.}{tab:rt-08}%
  {L{1.8cm} Y L{2cm} L{1.4cm} L{2.2cm} L{1.5cm} L{2cm}}%
  {\cab{Código} & \cab{Exigencia} & \cab{Carácter} & \cab{Cumple} & \cab{Componente} & \cab{Sección} & \cab{Evidencia}}
  RT-08.01 & El PROPONENTE especificará el equipamiento de cómputo con su marca, modelo de referencia, procesador, memoria, almacenamiento local, interfaces y consumo, junto con el cálculo de dimensionamiento que lo sustenta. & Obligatorio &  &  &  &  \\
  RT-08.02 & El almacenamiento será redundante, con tolerancia declarada a la falla de discos, control de errores y monitoreo predictivo de salud de los medios. & Obligatorio &  &  &  &  \\
  RT-08.03 & Los conmutadores de núcleo, los cortafuegos y los balanceadores de carga estarán en configuración de alta disponibilidad, sin punto único de falla. & Obligatorio &  &  &  &  \\
  RT-08.04 & Todo el equipamiento contará con fuentes de poder redundantes y conexión a circuitos eléctricos distintos. & Obligatorio &  &  &  &  \\
  RT-08.05 & El PROPONENTE declarará el margen de crecimiento del dimensionamiento propuesto, expresado como porcentaje sobre la carga proyectada del caso, y el procedimiento de ampliación. & Obligatorio &  &  &  &  \\
  RT-08.06 & El equipamiento será nuevo, sin uso previo, con garantía de fábrica vigente desde la recepción conforme. & Obligatorio &  &  &  &  \\
  RT-08.07 & El PROPONENTE especificará las estaciones de trabajo requeridas para la operación y la administración de la plataforma, en la cantidad que determine el dimensionamiento del caso, con monitores duales. & Obligatorio &  &  &  &  \\
  RT-08.08 & Los puestos de trabajo cumplirán las condiciones ergonómicas de la NCh 2527 y los equipos contarán con certificación de eficiencia energética. & Obligatorio &  &  &  &  \\
  RT-08.09 & Las estaciones estarán gestionadas de forma centralizada, con cifrado de disco, control de dispositivos extraíbles, antivirus con detección y respuesta y actualización automatizada. & Obligatorio &  &  &  &  \\
  RT-08.10 & El PROPONENTE especificará cada dispositivo de terreno con marca, modelo de referencia, cantidad, características mínimas, accesorios, consumibles y costo unitario estimado, aun cuando su compra sea de cargo del CLIENTE. & Obligatorio &  &  &  &  \\
  RT-08.11 & La especificación considerará las condiciones reales de uso del caso: intemperie, humedad, polvo, vibración, temperatura, uso con guantes, luminosidad y autonomía de batería requerida por turno. & Obligatorio &  &  &  &  \\
  RT-08.12 & Los dispositivos declararán su grado de protección contra polvo y agua y su resistencia a caídas, coherentes con el entorno de operación. & Obligatorio &  &  &  &  \\
  RT-08.13 & El PROPONENTE indicará el ciclo de vida esperado de cada dispositivo, la disponibilidad de repuestos y el plan de reposición durante los 56 meses del Contrato. & Obligatorio &  &  &  &  \\
  RT-08.14 & Los dispositivos se integrarán a la gestión centralizada de flota exigida en RT-03.18. & Obligatorio &  &  &  &  \\
  RT-08.15 & El PROPONENTE proveerá una unidad de cada tipo de dispositivo especificado para pruebas de aceptación por parte del CLIENTE, antes de la compra masiva. & Deseable &  &  &  &  \\
  RT-08.16 & El PROPONENTE presentará el plan de ciclo de vida del equipamiento: recepción, puesta en servicio, mantención, actualización, retiro y disposición final. & Obligatorio &  &  &  &  \\
  RT-08.17 & Todo medio de almacenamiento que salga de servicio será borrado de forma segura y verificable, con certificado de destrucción o de sanitización entregado al CLIENTE. & Obligatorio &  &  &  &  \\
  RT-08.18 & La disposición final de equipamiento electrónico se realizará con gestor autorizado, conforme a la normativa de residuos aplicable, con certificado de disposición. & Obligatorio &  &  &  &  \\
  RT-08.19 & Se valorará una estrategia de reacondicionamiento o de extensión de vida útil que reduzca el impacto ambiental, cuantificada en la propuesta. & Deseable &  &  &  &  \\
\end{tablalafrox}
Los requisitos del capítulo 09 se reproducen en la Tabla~\ref{tab:rt-09}.\par
\begin{tablalafrox}{Requisitos transversales, capítulo 09. Fuente: Bases Técnicas Transversales.}{tab:rt-09}%
  {L{1.8cm} Y L{2cm} L{1.4cm} L{2.2cm} L{1.5cm} L{2cm}}%
  {\cab{Código} & \cab{Exigencia} & \cab{Carácter} & \cab{Cumple} & \cab{Componente} & \cab{Sección} & \cab{Evidencia}}
  RT-09.01 & El PROPONENTE presentará el cálculo de capacidad que sustenta su dimensionamiento, con los supuestos de usuarios concurrentes, transacciones por segundo, volumen de datos y crecimiento anual, tomados de la volumetría del caso. & Obligatorio &  &  &  &  \\
  RT-09.02 & La solución soportará la concurrencia y el volumen de transacciones que fije el caso, y mantendrá los umbrales del numeral 9.1 bajo esa carga. & Según caso &  &  &  &  \\
  RT-09.03 & La solución soportará, sin rediseño, un crecimiento de al menos tres veces la volumetría inicial del caso en un horizonte de tres años. & Obligatorio &  &  &  &  \\
  RT-09.04 & El escalamiento de las capas de aplicación e integración será horizontal y automático, con tiempo de reacción declarado y sin pérdida de transacciones en curso. & Obligatorio &  &  &  &  \\
  RT-09.05 & El PROPONENTE identificará el componente que primero se convertirá en cuello de botella al crecer la carga y explicará cómo lo detectará y cómo lo resolverá. & Obligatorio &  &  &  &  \\
  RT-09.06 & Se ejecutarán pruebas de carga sobre Preproducción con un volumen equivalente a 1,5 veces el peak declarado, y pruebas de estrés hasta identificar el punto de quiebre de la solución. & Obligatorio &  &  &  &  \\
  RT-09.07 & El informe de pruebas de carga incluirá la curva de tiempo de respuesta frente a carga, el punto de saturación, el consumo de recursos y el comportamiento durante y después del peak. & Obligatorio &  &  &  &  \\
  RT-09.08 & La solución degradará de forma controlada al superarse la capacidad: encolamiento, limitación de tasa y mensaje explícito a la persona usuaria, nunca error genérico ni pérdida silenciosa de transacciones. & Obligatorio &  &  &  &  \\
  RT-09.09 & Se gestionará la capacidad durante la Operación con proyección trimestral de crecimiento, alertas anticipadas de agotamiento y propuesta de ajuste de dimensionamiento y de costo. & Obligatorio &  &  &  &  \\
  RT-09.10 & Se valorará la existencia de pruebas de carga automatizadas ejecutadas de forma periódica en el flujo de integración continua, con detección de regresiones de desempeño. & Deseable &  &  &  &  \\
\end{tablalafrox}
Los requisitos del capítulo 10 se reproducen en la Tabla~\ref{tab:rt-10}.\par
\begin{tablalafrox}{Requisitos transversales, capítulo 10. Fuente: Bases Técnicas Transversales.}{tab:rt-10}%
  {L{1.8cm} Y L{2cm} L{1.4cm} L{2.2cm} L{1.5cm} L{2cm}}%
  {\cab{Código} & \cab{Exigencia} & \cab{Carácter} & \cab{Cumple} & \cab{Componente} & \cab{Sección} & \cab{Evidencia}}
  RT-10.01 & La solución alcanzará una disponibilidad mensual mínima de 99,9 \% para los servicios clasificados como críticos, medida sobre la transacción de negocio de extremo a extremo. & Obligatorio &  &  &  &  \\
  RT-10.02 & El PROPONENTE clasificará cada servicio de la solución en crítico, alto, medio o bajo, justificando la clasificación con el impacto operacional de su indisponibilidad, y aplicará a cada uno el nivel de servicio correspondiente del Artículo 78° de las Bases Administrativas. & Obligatorio &  &  &  &  \\
  RT-10.03 & El plan de continuidad del negocio se elaborará conforme a ISO 22301, con análisis de impacto en el negocio, escenarios de contingencia, procedimientos manuales de respaldo y criterios de activación. & Obligatorio &  &  &  &  \\
  RT-10.04 & La continuidad TIC se estructurará conforme a ISO/IEC 27031, articulada con el plan de recuperación ante desastres del Capítulo 7. & Obligatorio &  &  &  &  \\
  RT-10.05 & Los mantenimientos programados se ejecutarán fuera de la ventana operacional crítica que defina el caso, con aviso previo mínimo de diez días hábiles. & Según caso &  &  &  &  \\
  RT-10.06 & La solución permitirá desplegar cambios sin interrupción del servicio. Las ventanas de indisponibilidad programada serán excepcionales y deberán justificarse caso a caso. & Obligatorio &  &  &  &  \\
  RT-10.07 & Se ejecutarán pruebas de resiliencia mediante inyección controlada de fallas ---caída de instancia, de zona, de dependencia externa, latencia elevada, saturación de disco--- antes de cada paso a producción y al menos una vez por semestre durante la Operación. & Obligatorio &  &  &  &  \\
  RT-10.08 & El PROPONENTE documentará, por cada dependencia externa, el comportamiento de la solución cuando esa dependencia no responde, responde con error o responde con lentitud. & Obligatorio &  &  &  &  \\
  RT-10.09 & Se declarará un presupuesto de error por servicio crítico y su vinculación con el ritmo de despliegue de cambios. & Deseable &  &  &  &  \\
\end{tablalafrox}
Los requisitos del capítulo 11 se reproducen en la Tabla~\ref{tab:rt-11}.\par
\begin{tablalafrox}{Requisitos transversales, capítulo 11. Fuente: Bases Técnicas Transversales.}{tab:rt-11}%
  {L{1.8cm} Y L{2cm} L{1.4cm} L{2.2cm} L{1.5cm} L{2cm}}%
  {\cab{Código} & \cab{Exigencia} & \cab{Carácter} & \cab{Cumple} & \cab{Componente} & \cab{Sección} & \cab{Evidencia}}
  RT-11.01 & La arquitectura de seguridad se basará en el modelo Zero Trust conforme a NIST SP 800-207: verificación explícita de cada solicitud, privilegio mínimo y presunción de compromiso. & Obligatorio &  &  &  &  \\
  RT-11.02 & El PROPONENTE entregará un modelado de amenazas documentado por cada componente y por cada integración externa, con metodología declarada (STRIDE u otra), y lo actualizará ante cada cambio arquitectónico relevante. & Obligatorio &  &  &  &  \\
  RT-11.03 & La información del CLIENTE se clasificará por nivel de sensibilidad, con controles diferenciados por nivel documentados en una matriz. & Obligatorio &  &  &  &  \\
  RT-11.04 & Existirá un programa de gestión de vulnerabilidades con escaneo continuo y plazos máximos de remediación de 7 días corridos para vulnerabilidades críticas, 15 días para altas y 30 días para medias, contados desde su publicación o detección. & Obligatorio &  &  &  &  \\
  RT-11.05 & El PROPONENTE mantendrá una matriz de controles de seguridad trazable a ISO/IEC 27001 e ISO/IEC 27002, indicando el control, su implementación concreta en la solución y la evidencia que lo acredita. & Obligatorio &  &  &  &  \\
  RT-11.06 & Se aplicarán los controles de ISO/IEC 27017 para los servicios en nube y de ISO/IEC 27018 para el tratamiento de datos personales en nube. & Obligatorio &  &  &  &  \\
  RT-11.07 & La publicación de servicios se realizará exclusivamente a través de la capa de borde, con red de distribución de contenidos, cortafuegos de aplicaciones web con reglas gestionadas y personalizadas, y protección contra denegación de servicio distribuida en capas 3, 4 y 7. & Obligatorio &  &  &  &  \\
  RT-11.08 & El cifrado en tránsito empleará TLS 1.3, con prohibición expresa de TLS 1.0 y 1.1, conjuntos de cifrado modernos, HSTS con precarga y gestión automatizada de certificados con rotación y alerta anticipada de vencimiento. & Obligatorio &  &  &  &  \\
  RT-11.09 & La totalidad de los datos en reposo estará cifrada, con claves gestionadas en un servicio de gestión de claves o en un módulo de seguridad de hardware, política de rotación declarada y separación de funciones en la custodia de claves. & Obligatorio &  &  &  &  \\
  RT-11.10 & Los datos de categoría sensible que el caso identifique se cifrarán adicionalmente a nivel de campo, de modo que el acceso a la base de datos no revele su contenido. & Según caso &  &  &  &  \\
  RT-11.11 & La puerta de enlace de servicios aplicará autenticación, autorización, cuotas, límites de tasa, validación de esquema e inspección de carga útil. & Obligatorio &  &  &  &  \\
  RT-11.12 & Los puntos de entrada públicos contarán con protección contra bots y abuso automatizado, con reto progresivo que no degrade la accesibilidad ni bloquee a personas usuarias legítimas. & Obligatorio &  &  &  &  \\
  RT-11.13 & El PROPONENTE declarará la superficie de exposición completa de la solución: cada nombre de dominio, puerto y servicio alcanzable desde fuera de la red del CLIENTE. & Obligatorio &  &  &  &  \\
  RT-11.14 & Los eventos de seguridad se registrarán de forma centralizada e inalterable, con retención mínima de doce meses en línea y veinticuatro meses adicionales en archivo recuperable. & Obligatorio &  &  &  &  \\
  RT-11.15 & Los eventos se correlacionarán en una plataforma SIEM, con casos de uso de detección definidos específicamente para el proceso de negocio del caso, y no sólo genéricos de infraestructura. & Obligatorio &  &  &  &  \\
  RT-11.16 & Se implementará detección y respuesta en puntos finales y en cargas de trabajo, tanto en nube como on-premise. & Obligatorio &  &  &  &  \\
  RT-11.17 & El PROPONENTE dispondrá de un centro de operaciones de seguridad con cobertura 24$\times$7, propio o subcontratado, y declarará su ubicación, dotación y procedimientos. & Obligatorio &  &  &  &  \\
  RT-11.18 & El plan de respuesta a incidentes de seguridad definirá clasificación, cadena de escalamiento, plazos, responsables y protocolo de comunicación al CLIENTE dentro de las dos horas de detectado un incidente de severidad crítica. & Obligatorio &  &  &  &  \\
  RT-11.19 & Toda brecha de seguridad o de datos personales se notificará al CLIENTE dentro de las 24 horas de su detección, con informe preliminar, y el análisis de causa raíz se entregará dentro de los cinco días hábiles siguientes. & Obligatorio &  &  &  &  \\
  RT-11.20 & Se ejecutarán pruebas de intrusión por un tercero independiente del ADJUDICATARIO, anualmente y antes de cada paso a producción, con entrega íntegra del informe al CLIENTE y plan de remediación con plazos. & Obligatorio &  &  &  &  \\
  RT-11.21 & Se realizarán ejercicios de simulación de incidente con participación del CLIENTE al menos una vez al año durante la Operación. & Deseable &  &  &  &  \\
  RT-11.22 & El flujo de integración continua incorporará análisis estático de código, análisis de composición de software, análisis dinámico y escaneo de imágenes de contenedor, con criterios de bloqueo automático del despliegue ante hallazgos críticos. & Obligatorio &  &  &  &  \\
  RT-11.23 & Cada versión liberada se acompañará de su inventario de componentes de software en formato CycloneDX o SPDX, entregado al CLIENTE. & Obligatorio &  &  &  &  \\
  RT-11.24 & Los artefactos se firmarán y su procedencia se verificará conforme a SLSA nivel 3 o superior. & Obligatorio &  &  &  &  \\
  RT-11.25 & Queda prohibido el uso de datos productivos reales en ambientes no productivos sin anonimización o seudonimización verificable. & Obligatorio &  &  &  &  \\
  RT-11.26 & El PROPONENTE declarará el proceso de aprobación de nuevas dependencias de terceros, incluyendo criterios de licencia, mantención activa y ausencia de vulnerabilidades conocidas. & Obligatorio &  &  &  &  \\
  RT-11.27 & Las personas desarrolladoras no tendrán acceso interactivo directo al ambiente de producción. Todo acceso excepcional será temporal, aprobado, registrado y con sesión grabada. & Obligatorio &  &  &  &  \\
  RT-11.28 & Se aplicará el marco OWASP SAMM o equivalente para medir y mejorar la madurez del proceso de desarrollo seguro, con evaluación inicial y reevaluación anual. & Deseable &  &  &  &  \\
\end{tablalafrox}
Los requisitos del capítulo 12 se reproducen en la Tabla~\ref{tab:rt-12}.\par
\begin{tablalafrox}{Requisitos transversales, capítulo 12. Fuente: Bases Técnicas Transversales.}{tab:rt-12}%
  {L{1.8cm} Y L{2cm} L{1.4cm} L{2.2cm} L{1.5cm} L{2cm}}%
  {\cab{Código} & \cab{Exigencia} & \cab{Carácter} & \cab{Cumple} & \cab{Componente} & \cab{Sección} & \cab{Evidencia}}
  RT-12.01 & La gestión de identidad será centralizada, con federación mediante OpenID Connect y OAuth 2.1, o SAML 2.0 cuando la integración con el CLIENTE lo requiera, e integración con el directorio corporativo del CLIENTE por LDAP o su equivalente en la nube. & Obligatorio &  &  &  &  \\
  RT-12.02 & Existirá inicio de sesión único para todos los módulos de la solución, con cierre de sesión propagado a todos ellos. & Obligatorio &  &  &  &  \\
  RT-12.03 & La autenticación multifactor será obligatoria para personas usuarias administradoras, para todo acceso privilegiado y para todo acceso originado fuera de la red corporativa. & Obligatorio &  &  &  &  \\
  RT-12.04 & Se soportarán factores resistentes a la suplantación de identidad, tipo FIDO2 o claves de acceso, al menos para los perfiles administradores. & Deseable &  &  &  &  \\
  RT-12.05 & El control de acceso será basado en roles, complementado con control basado en atributos donde el proceso lo exija, con matriz de segregación de funciones documentada y verificable. & Obligatorio &  &  &  &  \\
  RT-12.06 & Los accesos privilegiados se gestionarán con elevación temporal a demanda, aprobación previa y grabación de sesión para las operaciones de mayor riesgo. & Obligatorio &  &  &  &  \\
  RT-12.07 & La política de sesión declarará duración máxima, caducidad por inactividad, renovación de la credencial de sesión tras la autenticación, revocación inmediata y control de sesiones concurrentes. & Obligatorio &  &  &  &  \\
  RT-12.08 & Las credenciales de sesión serán firmadas y de vida breve, con credencial de refresco rotatoria. Queda prohibido transportar identificadores de sesión en la ruta de la dirección web. & Obligatorio &  &  &  &  \\
  RT-12.09 & Se registrará auditoría completa del ciclo de vida de la identidad: creación, modificación, elevación, bloqueo y baja de cuentas, con no repudio y retención declarada. & Obligatorio &  &  &  &  \\
  RT-12.10 & El aprovisionamiento y el desaprovisionamiento estarán automatizados y ligados al ciclo de vida laboral, con baja efectiva en un plazo no superior a 24 horas desde la desvinculación. & Obligatorio &  &  &  &  \\
  RT-12.11 & El mecanismo de autenticación se adecuará al perfil operacional real descrito en el caso: entornos de terreno, uso con guantes, baja alfabetización digital, dispositivos compartidos por turno y ausencia de correo electrónico personal. & Según caso &  &  &  &  \\
  RT-12.12 & Las personas usuarias externas del CLIENTE ---clientes, proveedores, pacientes, productores, según el caso--- dispondrán de un mecanismo de registro, verificación de identidad y recuperación de acceso autoservido y seguro. & Según caso &  &  &  &  \\
  RT-12.13 & El PROPONENTE declarará el procedimiento de acceso de emergencia (cuenta de último recurso), su custodia, su control y su auditoría. & Obligatorio &  &  &  &  \\
\end{tablalafrox}
Los requisitos del capítulo 13 se reproducen en la Tabla~\ref{tab:rt-13}.\par
\begin{tablalafrox}{Requisitos transversales, capítulo 13. Fuente: Bases Técnicas Transversales.}{tab:rt-13}%
  {L{1.8cm} Y L{2cm} L{1.4cm} L{2.2cm} L{1.5cm} L{2cm}}%
  {\cab{Código} & \cab{Exigencia} & \cab{Carácter} & \cab{Cumple} & \cab{Componente} & \cab{Sección} & \cab{Evidencia}}
  RT-13.01 & Todas las interfaces destinadas a personas usuarias cumplirán WCAG 2.2 nivel AA, verificado con herramientas automatizadas y con pruebas manuales, con informe de conformidad entregable. & Obligatorio &  &  &  &  \\
  RT-13.02 & El diseño será responsivo y funcionará correctamente en escritorio, tableta y teléfono, con puntos de quiebre coherentes y adaptación inteligente del contenido, no simple reducción. & Obligatorio &  &  &  &  \\
  RT-13.03 & El PROPONENTE ejecutará investigación con personas usuarias reales del CLIENTE, prototipado y pruebas de usabilidad antes de la construcción definitiva, y documentará los hallazgos y los cambios de diseño que produjeron. & Obligatorio &  &  &  &  \\
  RT-13.04 & Se comprometerán indicadores de usabilidad medibles: tiempo máximo de la transacción operacional crítica, número máximo de pasos, tasa de error tolerada y tiempo de aprendizaje esperado por perfil. & Obligatorio &  &  &  &  \\
  RT-13.05 & Ninguna funcionalidad principal requerirá más de tres interacciones desde la pantalla de inicio del perfil correspondiente. & Obligatorio &  &  &  &  \\
  RT-13.06 & La solución entregará retroalimentación visual clara ante cada acción, manejará los errores de forma comprensible ---indicando qué ocurrió y qué hacer--- y evitará mensajes técnicos dirigidos a la persona usuaria final. & Obligatorio &  &  &  &  \\
  RT-13.07 & La solución soportará a personas usuarias con baja alfabetización digital: alto contraste, objetivos táctiles de al menos 44 $\times$ 44 píxeles, iconografía acompañada de texto, y flujos guiados paso a paso. & Obligatorio &  &  &  &  \\
  RT-13.08 & Cuando el caso lo requiera, las interfaces de terreno operarán con guantes, a la intemperie, con luminosidad variable y sin conexión. & Según caso &  &  &  &  \\
  RT-13.09 & Existirá un sistema de diseño documentado con paleta de colores acotada, jerarquía tipográfica de a lo más dos familias, iconografía coherente, retícula y componentes reutilizables. & Obligatorio &  &  &  &  \\
  RT-13.10 & El PROPONENTE declarará la matriz de navegadores y de versiones soportadas y su política de actualización. & Obligatorio &  &  &  &  \\
  RT-13.11 & La solución será navegable íntegramente por teclado, con orden de foco lógico y atajos para las operaciones frecuentes. & Obligatorio &  &  &  &  \\
  RT-13.12 & Se valorará el soporte de modo oscuro, de personalización de la interfaz por persona usuaria y de múltiples idiomas cuando el caso lo justifique. & Deseable &  &  &  &  \\
\end{tablalafrox}
Los requisitos del capítulo 14 se reproducen en la Tabla~\ref{tab:rt-14}.\par
\begin{tablalafrox}{Requisitos transversales, capítulo 14. Fuente: Bases Técnicas Transversales.}{tab:rt-14}%
  {L{1.8cm} Y L{2cm} L{1.4cm} L{2.2cm} L{1.5cm} L{2cm}}%
  {\cab{Código} & \cab{Exigencia} & \cab{Carácter} & \cab{Cumple} & \cab{Componente} & \cab{Sección} & \cab{Evidencia}}
  RT-14.01 & La observabilidad será unificada para nube y on-premise, con métricas, registros y trazas distribuidas correlacionadas por un identificador único de transacción, instrumentadas conforme a OpenTelemetry. & Obligatorio &  &  &  &  \\
  RT-14.02 & El CLIENTE dispondrá de acceso propio y permanente a los tableros operacionales y de negocio, con datos en tiempo real y capacidad de exportación. & Obligatorio &  &  &  &  \\
  RT-14.03 & Los indicadores de nivel de servicio se medirán sobre la experiencia real de la persona usuaria y no sobre pruebas sintéticas, sin perjuicio de que estas se empleen como complemento. & Obligatorio &  &  &  &  \\
  RT-14.04 & El alertamiento se basará en síntomas de negocio y no sólo en umbrales de infraestructura, con supresión de ruido, agrupación, escalamiento automático y turnos de disponibilidad declarados. & Obligatorio &  &  &  &  \\
  RT-14.05 & Existirá un libro de operación y una guía de resolución documentados para cada escenario de falla previsible, con automatización progresiva de las tareas repetitivas. & Obligatorio &  &  &  &  \\
  RT-14.06 & Todo incidente crítico dará lugar a un análisis de causa raíz obligatorio, con informe entregable dentro de cinco días hábiles y seguimiento de las acciones correctivas hasta su cierre. & Obligatorio &  &  &  &  \\
  RT-14.07 & Los registros de la solución no contendrán datos personales sensibles ni credenciales, y su acceso estará controlado y auditado. & Obligatorio &  &  &  &  \\
  RT-14.08 & El PROPONENTE declarará la retención de métricas, registros y trazas, y su costo asociado, distinguiendo el almacenamiento en línea del archivado. & Obligatorio &  &  &  &  \\
  RT-14.09 & Se valorará la detección proactiva de anomalías mediante análisis del comportamiento histórico, con alerta antes de que el incidente afecte a la operación. & Deseable &  &  &  &  \\
\end{tablalafrox}
Los requisitos del capítulo 15 se reproducen en la Tabla~\ref{tab:rt-15}.\par
\begin{tablalafrox}{Requisitos transversales, capítulo 15. Fuente: Bases Técnicas Transversales.}{tab:rt-15}%
  {L{1.8cm} Y L{2cm} L{1.4cm} L{2.2cm} L{1.5cm} L{2cm}}%
  {\cab{Código} & \cab{Exigencia} & \cab{Carácter} & \cab{Cumple} & \cab{Componente} & \cab{Sección} & \cab{Evidencia}}
  RT-15.01 & El PROPONENTE dimensionará la infraestructura ajustada a la demanda real, evitando capacidad ociosa permanente, y declarará el factor de utilización proyectado. & Obligatorio &  &  &  &  \\
  RT-15.02 & Los ambientes no productivos se apagarán o reducirán fuera del horario de uso. & Obligatorio &  &  &  &  \\
  RT-15.03 & El PROPONENTE estimará la huella de carbono anual de la operación de la solución y declarará la metodología empleada. & Obligatorio &  &  &  &  \\
  RT-15.04 & Se declarará la eficiencia en el uso de la energía (PUE) del recinto on-premise y la intensidad de carbono de la región de nube escogida. & Obligatorio &  &  &  &  \\
  RT-15.05 & Se valorará la elección de regiones de nube con menor intensidad de carbono cuando la latencia y la regulación lo permitan, con el análisis comparativo correspondiente. & Deseable &  &  &  &  \\
  RT-15.06 & Se valorará la definición de metas de reducción del consumo durante la Operación, con medición y reporte anual. & Deseable &  &  &  &  \\
  RT-15.07 & Las certificaciones se acreditarán con copia del certificado vigente y con el código de verificación del organismo emisor cuando exista. & Obligatorio &  &  &  &  \\
  RT-15.08 & Las personas certificadas formarán parte del equipo efectivamente asignado al PROYECTO, con dedicación declarada. No se aceptará acreditar personal que no participe. & Obligatorio &  &  &  &  \\
  RT-15.09 & El ADJUDICATARIO mantendrá vigentes estas certificaciones durante todo el período contractual y las repondrá ante la salida de una persona certificada, conforme al Artículo 76° de las Bases Administrativas. & Obligatorio &  &  &  &  \\
\end{tablalafrox}
Los requisitos del capítulo 16 se reproducen en la Tabla~\ref{tab:rt-16}.\par
\begin{tablalafrox}{Requisitos transversales, capítulo 16. Fuente: Bases Técnicas Transversales.}{tab:rt-16}%
  {L{1.8cm} Y L{2cm} L{1.4cm} L{2.2cm} L{1.5cm} L{2cm}}%
  {\cab{Código} & \cab{Exigencia} & \cab{Carácter} & \cab{Cumple} & \cab{Componente} & \cab{Sección} & \cab{Evidencia}}
  RT-16.01 & La solución dispondrá de un módulo de administración que permita al CLIENTE, sin intervención del ADJUDICATARIO, gestionar personas usuarias, roles, permisos, unidades organizacionales y sus jerarquías. & Obligatorio &  &  &  &  \\
  RT-16.02 & Las reglas de negocio parametrizables ---umbrales, plazos, montos, tolerancias, catálogos, listas de valores, textos de notificación--- serán configurables desde la interfaz de administración, con control de versiones y registro de quién cambió qué y cuándo. & Obligatorio &  &  &  &  \\
  RT-16.03 & Todo cambio de parámetro con impacto operacional requerirá aprobación de un segundo perfil y quedará registrado con su justificación. & Obligatorio &  &  &  &  \\
  RT-16.04 & El PROPONENTE declarará expresamente qué elementos son parametrizables y cuáles requieren desarrollo. Presentar como parametrizable lo que exige desarrollo se evaluará como observación grave. & Obligatorio &  &  &  &  \\
  RT-16.05 & Existirá un ambiente de simulación que permita probar el efecto de un cambio de parámetro antes de aplicarlo a producción. & Deseable &  &  &  &  \\
  RT-16.06 & Toda operación que cree, modifique o elimine información quedará registrada con identificación de la persona o del sistema que la ejecutó, fecha y hora con zona horaria, origen, valores anteriores y valores posteriores. & Obligatorio &  &  &  &  \\
  RT-16.07 & El registro de auditoría será inalterable y no podrá ser modificado ni eliminado por ningún perfil, incluido el administrador de la plataforma. & Obligatorio &  &  &  &  \\
  RT-16.08 & El CLIENTE podrá consultar y exportar la auditoría desde la interfaz, con filtros por persona, período, entidad y tipo de operación, sin requerir acceso a la base de datos. & Obligatorio &  &  &  &  \\
  RT-16.09 & Las consultas a información sensible quedarán registradas, no sólo las modificaciones. & Según caso &  &  &  &  \\
  RT-16.10 & El período de retención de la auditoría será el que fije el caso y, en su defecto, no inferior a cinco años. & Según caso &  &  &  &  \\
  RT-16.11 & La solución soportará flujos de trabajo con estados, transiciones, responsables, plazos, escalamiento automático por vencimiento y delegación por ausencia. & Obligatorio &  &  &  &  \\
  RT-16.12 & Los flujos serán configurables por el CLIENTE sin desarrollo, al menos en lo relativo a responsables, plazos y niveles de aprobación. & Obligatorio &  &  &  &  \\
  RT-16.13 & Toda solicitud pendiente será visible para su responsable en una bandeja de tareas unificada, con priorización y alerta de vencimiento. & Obligatorio &  &  &  &  \\
  RT-16.14 & El motor de reglas permitirá definir condiciones de negocio evaluables sin recompilación, con trazabilidad de qué regla se aplicó a cada transacción. & Deseable &  &  &  &  \\
  RT-16.15 & La solución gestionará documentos con versionado, metadatos, control de acceso, búsqueda por contenido y por metadato, y previsualización sin descarga. & Obligatorio &  &  &  &  \\
  RT-16.16 & Los documentos se almacenarán cifrados, con verificación de integridad y retención conforme a la política del caso. & Obligatorio &  &  &  &  \\
  RT-16.17 & La solución soportará firma electrónica conforme a la Ley N° 19.799, con firma avanzada para los actos que el caso lo requiera, y verificación de validez del certificado al momento de la firma. & Según caso &  &  &  &  \\
  RT-16.18 & Se generará el sello de tiempo y se conservará la evidencia de firma que permita verificar el documento con posterioridad al vencimiento del certificado. & Según caso &  &  &  &  \\
  RT-16.19 & La solución generará documentos a partir de plantillas administrables por el CLIENTE, con datos de la transacción y salida en formato abierto. & Obligatorio &  &  &  &  \\
  RT-16.20 & La solución enviará notificaciones por al menos tres canales: correo electrónico, notificación en la aplicación y mensajería instantánea o SMS, según lo que el caso requiera. & Obligatorio &  &  &  &  \\
  RT-16.21 & Las plantillas de notificación serán administrables por el CLIENTE, con variables de la transacción, y versionadas. & Obligatorio &  &  &  &  \\
  RT-16.22 & Cada persona usuaria podrá configurar sus preferencias de canal y de frecuencia, respetando las notificaciones que el CLIENTE defina como obligatorias. & Obligatorio &  &  &  &  \\
  RT-16.23 & El envío será asíncrono, con reintento ante falla, control de duplicados y registro de entrega, apertura y error por cada mensaje. & Obligatorio &  &  &  &  \\
  RT-16.24 & El PROPONENTE declarará el proveedor de cada canal, su costo unitario, su volumen proyectado y el tratamiento del costo variable en la Oferta Económica. & Obligatorio &  &  &  &  \\
  RT-16.25 & Las notificaciones respetarán la normativa de comunicaciones comerciales y permitirán la baja cuando corresponda. & Obligatorio &  &  &  &  \\
  RT-16.26 & Se valorará la integración con canales conversacionales que permitan a la persona usuaria responder y ejecutar acciones desde el propio canal. & Deseable &  &  &  &  \\
  RT-16.27 & La solución dispondrá de búsqueda global con indexación de texto completo, tolerancia a errores de escritura, filtros facetados y respeto del control de acceso de la persona que busca. & Obligatorio &  &  &  &  \\
  RT-16.28 & Los listados serán ordenables, filtrables, paginados y exportables en formatos abiertos, con el filtro aplicado reflejado en la exportación. & Obligatorio &  &  &  &  \\
  RT-16.29 & Las exportaciones de gran volumen se procesarán de forma asíncrona, con notificación al completarse y sin bloquear la sesión. & Obligatorio &  &  &  &  \\
  RT-16.30 & Toda exportación de información sensible quedará registrada en la auditoría, con identificación de quién exportó qué y cuándo. & Obligatorio &  &  &  &  \\
  RT-16.31 & La solución dispondrá de un portal público con la información que el caso determine, accesible sin autenticación, con los mismos estándares de accesibilidad y desempeño que el resto de la plataforma. & Según caso &  &  &  &  \\
  RT-16.32 & Las personas usuarias externas dispondrán de autoatención para las consultas de mayor frecuencia, evitando el contacto telefónico para operaciones simples. & Obligatorio &  &  &  &  \\
  RT-16.33 & El PROPONENTE estimará la reducción esperada del volumen de atención asistida por efecto de la autoatención y comprometerá el indicador. & Obligatorio &  &  &  &  \\
  RT-16.34 & El portal público resistirá picos de tráfico sin degradar los servicios transaccionales internos, mediante aislamiento de recursos y caché. & Obligatorio &  &  &  &  \\
\end{tablalafrox}
Los requisitos del capítulo 17 se reproducen en la Tabla~\ref{tab:rt-17}.\par
\begin{tablalafrox}{Requisitos transversales, capítulo 17. Fuente: Bases Técnicas Transversales.}{tab:rt-17}%
  {L{1.8cm} Y L{2cm} L{1.4cm} L{2.2cm} L{1.5cm} L{2cm}}%
  {\cab{Código} & \cab{Exigencia} & \cab{Carácter} & \cab{Cumple} & \cab{Componente} & \cab{Sección} & \cab{Evidencia}}
  RT-17.01 & La solución proveerá una aplicación móvil para los perfiles operacionales que el caso identifique, con funcionamiento sin conexión y sincronización diferida cuando la operación en terreno lo requiera. & Según caso &  &  &  &  \\
  RT-17.02 & El PROPONENTE declarará si la aplicación es nativa, híbrida o web progresiva, y justificará la decisión frente al requisito de operación desconectada, al acceso a periféricos y al costo de mantención. & Obligatorio &  &  &  &  \\
  RT-17.03 & La aplicación soportará las versiones de sistema operativo móvil vigentes y las dos anteriores, con política de actualización declarada. & Obligatorio &  &  &  &  \\
  RT-17.04 & La aplicación se distribuirá por las tiendas oficiales o mediante gestión de flota corporativa, con firma de la aplicación y verificación de integridad. & Obligatorio &  &  &  &  \\
  RT-17.05 & La información almacenada en el dispositivo estará cifrada, con borrado remoto y bloqueo ante pérdida o desvinculación de la persona usuaria. & Obligatorio &  &  &  &  \\
  RT-17.06 & La aplicación integrará los periféricos que el caso requiera: cámara, lector de códigos, NFC, GPS, impresora de etiquetas, balanza o báscula. & Según caso &  &  &  &  \\
  RT-17.07 & El consumo de datos móviles y de batería se optimizará y se declarará el consumo estimado por turno de trabajo. & Obligatorio &  &  &  &  \\
  RT-17.08 & Se valorará la disponibilidad de una versión de la interfaz para dispositivos de bajo costo o de generaciones anteriores, ampliando la cobertura de personas usuarias. & Deseable &  &  &  &  \\
\end{tablalafrox}
Los requisitos del capítulo 18 se reproducen en la Tabla~\ref{tab:rt-18}.\par
\begin{tablalafrox}{Requisitos transversales, capítulo 18. Fuente: Bases Técnicas Transversales.}{tab:rt-18}%
  {L{1.8cm} Y L{2cm} L{1.4cm} L{2.2cm} L{1.5cm} L{2cm}}%
  {\cab{Código} & \cab{Exigencia} & \cab{Carácter} & \cab{Cumple} & \cab{Componente} & \cab{Sección} & \cab{Evidencia}}
  RT-18.01 & El PROPONENTE declarará cada componente de inteligencia artificial: propósito, modelo o servicio empleado, proveedor, versión y ubicación de procesamiento de los datos. & Obligatorio &  &  &  &  \\
  RT-18.02 & Se garantizará contractualmente que los datos del CLIENTE no serán utilizados para entrenar modelos de terceros, salvo autorización expresa y escrita. & Obligatorio &  &  &  &  \\
  RT-18.03 & Se documentarán los límites de uso, los casos en que el resultado requiere validación humana previa y el procedimiento de supervisión. & Obligatorio &  &  &  &  \\
  RT-18.04 & Los riesgos de sesgo, alucinación, fuga de información y uso indebido se evaluarán y mitigarán conforme al NIST AI Risk Management Framework 1.0 y a la norma ISO/IEC 42001. & Obligatorio &  &  &  &  \\
  RT-18.05 & Toda interacción relevante con un componente de inteligencia artificial quedará registrada para efectos de auditoría y trazabilidad, incluyendo la entrada, la salida y la decisión humana posterior. & Obligatorio &  &  &  &  \\
  RT-18.06 & El componente podrá desactivarse sin comprometer la operación del resto de la solución, y existirá un procedimiento manual de respaldo para la función que automatiza. & Obligatorio &  &  &  &  \\
  RT-18.07 & Cuando el resultado se presente a una persona usuaria, se indicará expresamente que fue generado o sugerido de forma automática, y su nivel de confianza cuando el modelo lo provea. & Obligatorio &  &  &  &  \\
  RT-18.08 & Los modelos predictivos declararán sus variables de entrada, su métrica de desempeño, su línea base y su plan de reentrenamiento y de detección de deriva. & Obligatorio &  &  &  &  \\
  RT-18.09 & La responsabilidad por los resultados de los componentes de inteligencia artificial recae íntegramente en el ADJUDICATARIO, conforme al Artículo 86° de las Bases Administrativas. & Obligatorio &  &  &  &  \\
  RT-18.10 & Se valorará el uso de automatización robótica de procesos o de agentes para tareas repetitivas de back office, con el ahorro de horas cuantificado y reflejado en la Oferta Económica. & Deseable &  &  &  &  \\
\end{tablalafrox}
Los requisitos del capítulo 19 se reproducen en la Tabla~\ref{tab:rt-19}.\par
\begin{tablalafrox}{Requisitos transversales, capítulo 19. Fuente: Bases Técnicas Transversales.}{tab:rt-19}%
  {L{1.8cm} Y L{2cm} L{1.4cm} L{2.2cm} L{1.5cm} L{2cm}}%
  {\cab{Código} & \cab{Exigencia} & \cab{Carácter} & \cab{Cumple} & \cab{Componente} & \cab{Sección} & \cab{Evidencia}}
  RT-19.01 & El ADJUDICATARIO constituirá una oficina de gestión del proyecto con metodología declarada, basada en el PMBOK del Project Management Institute e integrando prácticas ágiles donde el PROPONENTE lo justifique. & Obligatorio &  &  &  &  \\
  RT-19.02 & El Jefe de Proyecto tendrá dedicación exclusiva durante toda la fase de implementación y facultades para comprometer al ADJUDICATARIO en materias de ejecución. & Obligatorio &  &  &  &  \\
  RT-19.03 & Existirá un procedimiento formal de control de cambios conforme al Artículo 72° de las Bases Administrativas, con registro de cambios, análisis de impacto y aprobación previa a la ejecución. & Obligatorio &  &  &  &  \\
  RT-19.04 & La gestión del riesgo seguirá la norma ISO 31000, con registro de riesgos vivo, revisión en cada Comité de Proyecto y planes de mitigación con responsable, plazo y disparador. & Obligatorio &  &  &  &  \\
  RT-19.05 & El ADJUDICATARIO habilitará un espacio colaborativo accesible al CLIENTE, con la documentación del proyecto, los entregables, las actas, el registro de riesgos y el registro de cambios siempre actualizados. & Obligatorio &  &  &  &  \\
  RT-19.06 & El ADJUDICATARIO entregará un informe mensual de avance con estado del cronograma, avance físico y financiero, entregables del período, desviaciones, riesgos, incidencias y compromisos del período siguiente. & Obligatorio &  &  &  &  \\
  RT-19.07 & El avance se medirá con valor ganado, reportando el índice de desempeño del cronograma y del costo, y no mediante declaración cualitativa de porcentaje de avance. & Obligatorio &  &  &  &  \\
  RT-19.08 & El CLIENTE dispondrá de un tablero de estado del proyecto actualizado, accesible en cualquier momento. & Obligatorio &  &  &  &  \\
  RT-19.09 & Las actas de todos los comités del Artículo 71° de las Bases Administrativas se levantarán dentro de los dos días hábiles siguientes y registrarán acuerdos, responsables y plazos. & Obligatorio &  &  &  &  \\
  RT-19.10 & Toda desviación superior al 10 \% en un hito se comunicará dentro de los cinco días hábiles de detectada, con plan de recuperación. & Obligatorio &  &  &  &  \\
\end{tablalafrox}
Los requisitos del capítulo 20 se reproducen en la Tabla~\ref{tab:rt-20}.\par
\begin{tablalafrox}{Requisitos transversales, capítulo 20. Fuente: Bases Técnicas Transversales.}{tab:rt-20}%
  {L{1.8cm} Y L{2cm} L{1.4cm} L{2.2cm} L{1.5cm} L{2cm}}%
  {\cab{Código} & \cab{Exigencia} & \cab{Carácter} & \cab{Cumple} & \cab{Componente} & \cab{Sección} & \cab{Evidencia}}
  RT-20.01 & El PROPONENTE presentará un plan de implantación por sitio y por perfil, coherente con el cronograma obligatorio del Artículo 17° de las Bases Administrativas. & Obligatorio &  &  &  &  \\
  RT-20.02 & La estrategia de paso a producción será gradual y reversible. Se declarará el criterio de avance entre olas y el procedimiento de reversión, con su tiempo de ejecución. & Obligatorio &  &  &  &  \\
  RT-20.03 & Durante la marcha blanca, la solución convivirá con la operación vigente del CLIENTE, con conciliación diaria y sin doble digitación no declarada. & Obligatorio &  &  &  &  \\
  RT-20.04 & El PROPONENTE definirá los indicadores diarios que se medirán durante la marcha blanca y sus umbrales de avance, conforme al Artículo 17.3 de las Bases Administrativas. & Obligatorio &  &  &  &  \\
  RT-20.05 & Se dispondrá de acompañamiento en terreno durante las primeras semanas de cada ola, con dotación declarada y decreciente según la curva de adopción. & Obligatorio &  &  &  &  \\
  RT-20.06 & Se establecerá un período de estabilización con atención reforzada tras cada paso a producción, con dotación y duración declaradas y sin costo adicional. & Obligatorio &  &  &  &  \\
  RT-20.07 & Existirá una definición de terminado acordada con el CLIENTE, aplicable a cada entregable, que incluya código, pruebas, documentación, seguridad y despliegue. & Obligatorio &  &  &  &  \\
  RT-20.08 & El protocolo de aceptación de cada hito se formalizará conforme al Formulario T-17, con criterios objetivos y verificables. & Obligatorio &  &  &  &  \\
\end{tablalafrox}
Los requisitos del capítulo 21 se reproducen en la Tabla~\ref{tab:rt-21}.\par
\begin{tablalafrox}{Requisitos transversales, capítulo 21. Fuente: Bases Técnicas Transversales.}{tab:rt-21}%
  {L{1.8cm} Y L{2cm} L{1.4cm} L{2.2cm} L{1.5cm} L{2cm}}%
  {\cab{Código} & \cab{Exigencia} & \cab{Carácter} & \cab{Cumple} & \cab{Componente} & \cab{Sección} & \cab{Evidencia}}
  RT-21.01 & El ADJUDICATARIO dispondrá de un centro de operaciones de red con cobertura 24$\times$7$\times$365, propio o subcontratado, y declarará su ubicación, dotación por turno y procedimientos. & Obligatorio &  &  &  &  \\
  RT-21.02 & Se designará un gerente de servicio dedicado como contraparte permanente del CLIENTE durante la fase de Operación. & Obligatorio &  &  &  &  \\
  RT-21.03 & El equipo de operación contará con especialistas por tecnología, nominados y con dedicación declarada. & Obligatorio &  &  &  &  \\
  RT-21.04 & Los procedimientos de operación estarán documentados y serán ejecutables por el personal del CLIENTE tras la transferencia de conocimiento. & Obligatorio &  &  &  &  \\
  RT-21.05 & El PROPONENTE dispondrá de un centro de atención adecuado para soportar integralmente las consultas, propio o subcontratado con un operador de nivel acreditado. & Obligatorio &  &  &  &  \\
  RT-21.06 & El centro de atención cumplirá un tiempo de atención al usuario final de 80 \% antes de 20 segundos y una resolución al primer contacto de al menos 70 \%. La tasa de abandono no superará el 5 \%. & Obligatorio &  &  &  &  \\
  RT-21.07 & El horario mínimo de atención será de 8:00 a 20:00 en días hábiles, ampliado a 24$\times$7 para los incidentes de severidad crítica y para la ventana operacional que defina el caso. & Según caso &  &  &  &  \\
  RT-21.08 & La atención cubrirá orientación funcional de distintos grados de complejidad, desde preguntas simples hasta situaciones complejas, y las preguntas técnicas más frecuentes sobre la aplicación y su entorno de uso. & Obligatorio &  &  &  &  \\
  RT-21.09 & Se mantendrá un registro histórico de las actividades de soporte que permita monitorear y gestionar los principales requerimientos, con análisis de tendencia mensual. & Obligatorio &  &  &  &  \\
  RT-21.10 & Existirá un proceso de aprendizaje del centro de soporte que acumule conocimiento sobre las complejidades de la operación y las necesidades de las personas usuarias, reflejado en la base de conocimiento. & Obligatorio &  &  &  &  \\
  RT-21.11 & Se proveerán servicios de capacitación en línea en modalidad de autoformación, con registro de avance para efectos de monitoreo. & Obligatorio &  &  &  &  \\
  RT-21.12 & Se habilitarán espacios de interacción entre las entidades y personas usuarias del sistema, que permitan compartir experiencias de uso. & Deseable &  &  &  &  \\
  RT-21.13 & Los indicadores clave del proceso de soporte estarán abiertos al CLIENTE, y los indicadores específicos serán accesibles según nivel de autorización. & Obligatorio &  &  &  &  \\
  RT-21.14 & El PROPONENTE dimensionará la dotación del centro de atención con fundamento cuantitativo, empleando teoría de colas o el modelo Erlang C, a partir del volumen de contactos proyectado del caso. & Obligatorio &  &  &  &  \\
  RT-21.15 & Existirá un canal único de registro de incidentes y solicitudes, con número de ticket, clasificación por severidad y seguimiento del ciclo de vida completo hasta el cierre conforme. & Obligatorio &  &  &  &  \\
  RT-21.16 & Cuando el caso comprenda sitios alejados, los especialistas de niveles 2 y 3 deberán trasladarse cuando la resolución lo requiera, por el medio más rápido disponible y con disponibilidad de tiempo suficiente, sin alterar la atención normal del resto de los sitios. El costo del traslado está incluido en la oferta. & Según caso &  &  &  &  \\
  RT-21.17 & El cierre de un ticket requerirá confirmación de la persona usuaria o transcurso del plazo de confirmación automática declarado. & Obligatorio &  &  &  &  \\
  RT-21.18 & El ADJUDICATARIO reportará mensualmente el cumplimiento de los niveles de servicio de atención, con el detalle por severidad y el análisis de los incumplimientos. & Obligatorio &  &  &  &  \\
  RT-21.19 & El PROPONENTE declarará el tamaño de la bolsa anual de horas de mantención evolutiva, su composición por perfil y el procedimiento de solicitud, estimación, aprobación y liquidación. & Obligatorio &  &  &  &  \\
  RT-21.20 & La mantención evolutiva se someterá al mismo estándar de calidad, pruebas y seguridad que el desarrollo original. & Obligatorio &  &  &  &  \\
  RT-21.21 & El PROPONENTE mantendrá un registro de deuda técnica cuantificado y destinará una fracción declarada de la capacidad de la Operación a reducirla. & Obligatorio &  &  &  &  \\
  RT-21.22 & Las actualizaciones de versión de los componentes de base se planificarán anualmente, con ventana acordada y plan de reversión. & Obligatorio &  &  &  &  \\
\end{tablalafrox}
Los requisitos del capítulo 22 se reproducen en la Tabla~\ref{tab:rt-22}.\par
\begin{tablalafrox}{Requisitos transversales, capítulo 22. Fuente: Bases Técnicas Transversales.}{tab:rt-22}%
  {L{1.8cm} Y L{2cm} L{1.4cm} L{2.2cm} L{1.5cm} L{2cm}}%
  {\cab{Código} & \cab{Exigencia} & \cab{Carácter} & \cab{Cumple} & \cab{Componente} & \cab{Sección} & \cab{Evidencia}}
  RT-22.01 & El plan de capacitación se estructurará por perfil: personas usuarias finales, usuarias avanzadas, administradoras, equipo técnico y soporte de niveles 1 y 2. & Obligatorio &  &  &  &  \\
  RT-22.02 & Las modalidades incluirán capacitación presencial en cada sitio de operación, sesiones en línea sincrónicas, autoformación y acompañamiento en puesto de trabajo durante la marcha blanca. & Obligatorio &  &  &  &  \\
  RT-22.03 & Todo el material de capacitación se entregará en español, en formato editable y de propiedad del CLIENTE: manuales por perfil, guías rápidas, preguntas frecuentes, videos tutoriales y base de conocimiento consultable. & Obligatorio &  &  &  &  \\
  RT-22.04 & La capacitación no podrá afectar la operación del CLIENTE: se programará por turnos y en horarios acordados, considerando la estacionalidad y la ventana operacional del caso. & Según caso &  &  &  &  \\
  RT-22.05 & Las personas usuarias administradoras y el equipo técnico del CLIENTE serán evaluados y certificados como condición para el cierre de cada marcha blanca. & Obligatorio &  &  &  &  \\
  RT-22.06 & Durante la Operación se ejecutarán al menos dos jornadas anuales de actualización y se capacitará al personal nuevo del CLIENTE sin costo adicional. & Obligatorio &  &  &  &  \\
  RT-22.07 & El ADJUDICATARIO proveerá acompañamiento y mentoría al equipo técnico del CLIENTE durante al menos seis meses posteriores al paso a producción de la Etapa 2. & Obligatorio &  &  &  &  \\
  RT-22.08 & La base de conocimiento será mantenida y actualizada durante todo el Contrato, con métricas de uso y de utilidad percibida. & Obligatorio &  &  &  &  \\
  RT-22.09 & Se valorará la existencia de un ambiente permanente de entrenamiento, con datos ficticios, disponible para la práctica sin riesgo. & Deseable &  &  &  &  \\
\end{tablalafrox}
Los requisitos del capítulo 23 se reproducen en la Tabla~\ref{tab:rt-23}.\par
\begin{tablalafrox}{Requisitos transversales, capítulo 23. Fuente: Bases Técnicas Transversales.}{tab:rt-23}%
  {L{1.8cm} Y L{2cm} L{1.4cm} L{2.2cm} L{1.5cm} L{2cm}}%
  {\cab{Código} & \cab{Exigencia} & \cab{Carácter} & \cab{Cumple} & \cab{Componente} & \cab{Sección} & \cab{Evidencia}}
  RT-23.01 & Información institucional: descripción detallada del giro principal, historia y trayectoria de la empresa con al menos tres años de antigüedad comprobable, misión, visión y valores, certificaciones y acreditaciones vigentes, y presencia geográfica y oficinas. & Obligatorio &  &  &  &  \\
  RT-23.02 & Experiencia y casos de éxito: portafolio de proyectos similares de los últimos cinco años, casos documentados con métricas verificables, testimonios de clientes con autorización de publicación, industrias atendidas con énfasis en el caso asignado, y volumen de transacciones o de personas usuarias gestionadas. & Obligatorio &  &  &  &  \\
  RT-23.03 & Equipo profesional: organigrama del equipo directivo, perfiles de socios y directores, currículos resumidos del equipo técnico clave y certificaciones profesionales del personal. & Obligatorio &  &  &  &  \\
  RT-23.04 & Capacidades técnicas: servicios y soluciones ofrecidas, conjunto tecnológico dominado, alianzas tecnológicas con proveedores de nube y de plataforma, metodologías de trabajo certificadas, e infraestructura y capacidad instalada. & Obligatorio &  &  &  &  \\
  RT-23.05 & El sitio será accesible conforme a WCAG 2.2 nivel AA, responsivo y con certificado TLS válido y vigente. & Obligatorio &  &  &  &  \\
  RT-23.06 & El sitio se mantendrá activo y disponible durante todo el proceso de licitación, desde el registro del participante hasta la adjudicación. & Obligatorio &  &  &  &  \\
  RT-23.07 & Material educativo digital: seminarios en línea, artículos técnicos o centro de recursos y documentación. & Deseable &  &  &  &  \\
  RT-23.08 & Demostración en línea, recorrido virtual de la solución o portal de soporte para clientes. & Deseable &  &  &  &  \\
  RT-23.09 & Métricas de disponibilidad y desempeño de los servicios del proponente publicadas en tiempo real. & Deseable &  &  &  &  \\
  RT-23.10 & Calculadora de retorno de la inversión o herramienta de estimación pertinente a la industria del caso. & Deseable &  &  &  &  \\
\end{tablalafrox}
Los requisitos del capítulo 24 se reproducen en la Tabla~\ref{tab:rt-24}.\par
\begin{tablalafrox}{Requisitos transversales, capítulo 24. Fuente: Bases Técnicas Transversales.}{tab:rt-24}%
  {L{1.8cm} Y L{2cm} L{1.4cm} L{2.2cm} L{1.5cm} L{2cm}}%
  {\cab{Código} & \cab{Exigencia} & \cab{Carácter} & \cab{Cumple} & \cab{Componente} & \cab{Sección} & \cab{Evidencia}}
  RT-24.01 & Todos los integrantes clave del equipo deberán aparecer individualmente presentándose, en toma de medio cuerpo, mirando directamente a la cámara, con vestimenta formal o de negocio informal, y con una duración mínima de aparición de 10 segundos por persona. & Obligatorio &  &  &  &  \\
  RT-24.02 & Cada vez que aparezca un integrante deberá mostrarse en pantalla su nombre completo, su cargo en el proyecto, sus años de experiencia y su especialización relevante. & Obligatorio &  &  &  &  \\
  RT-24.03 & Las certificaciones principales de cada integrante podrán mostrarse en pantalla. & Deseable &  &  &  &  \\
  RT-24.04 & El fondo será el mismo para todas las tomas de personas, con iluminación consistente y encuadre similar para todos los participantes. & Obligatorio &  &  &  &  \\
  RT-24.05 & La paleta de colores corporativa, la tipografía y el uso de la marca serán consistentes en todo el video, con logotipo visible en todas las escenas y datos de contacto en el encabezado o pie. & Obligatorio &  &  &  &  \\
  RT-24.06 & Los niveles de volumen estarán normalizados, con la misma calidad de grabación y sin variaciones bruscas de sonido. & Obligatorio &  &  &  &  \\
\end{tablalafrox}
Los requisitos del capítulo 25 se reproducen en la Tabla~\ref{tab:rt-25}.\par
\begin{tablalafrox}{Requisitos transversales, capítulo 25. Fuente: Bases Técnicas Transversales.}{tab:rt-25}%
  {L{1.8cm} Y L{2cm} L{1.4cm} L{2.2cm} L{1.5cm} L{2cm}}%
  {\cab{Código} & \cab{Exigencia} & \cab{Carácter} & \cab{Cumple} & \cab{Componente} & \cab{Sección} & \cab{Evidencia}}
  RT-25.01 & Navegación intuitiva sin necesidad de manual, con un máximo de tres interacciones para alcanzar cualquier funcionalidad principal. & Obligatorio &  &  &  &  \\
  RT-25.02 & Retroalimentación visual clara ante cada acción, estados de carga y transición fluidos, y manejo elegante de los errores. & Obligatorio &  &  &  &  \\
  RT-25.03 & Cumplimiento de WCAG 2.2 nivel AA: contraste adecuado, tamaños de fuente legibles, objetivos táctiles de al menos 44 $\times$ 44 píxeles, navegación completa por teclado, y textos alternativos y etiquetas de accesibilidad. & Obligatorio &  &  &  &  \\
  RT-25.04 & Diseño responsivo demostrado en escritorio (1920 $\times$ 1080 y 1366 $\times$ 768), tableta (768 $\times$ 1024 en vertical y horizontal) y teléfono (375 $\times$ 812), con puntos de quiebre coherentes. & Obligatorio &  &  &  &  \\
  RT-25.05 & Sistema de diseño documentado: paleta de a lo más cinco colores principales, a lo más dos familias tipográficas jerarquizadas, iconografía coherente, retícula y espaciado consistentes, y componentes reutilizables. & Obligatorio &  &  &  &  \\
  RT-25.06 & El prototipo se entregará como enlace navegable en línea, sin instalación ni complementos, compatible con los navegadores modernos vigentes. & Obligatorio &  &  &  &  \\
  RT-25.07 & El acceso al prototipo estará garantizado por un mínimo de seis meses desde su entrega. & Obligatorio &  &  &  &  \\
  RT-25.08 & El prototipo permitirá navegación completa entre pantallas, simulación de ingreso de datos, transiciones y animaciones básicas, estados de posado y activo, simulación de carga de archivos, elementos desplegables funcionales y simulación de validaciones. & Obligatorio &  &  &  &  \\
  RT-25.09 & Se entregará la arquitectura de información del prototipo: mapa del sitio completo, diagrama de navegación, taxonomía y nomenclatura, estructura de menús y jerarquía de la información. & Obligatorio &  &  &  &  \\
  RT-25.10 & Se valorará la posibilidad de exportar el prototipo a un documento interactivo para revisión sin conexión. & Deseable &  &  &  &  \\
\end{tablalafrox}
Los requisitos del capítulo 26 se reproducen en la Tabla~\ref{tab:rt-26}.\par
\begin{tablalafrox}{Requisitos transversales, capítulo 26. Fuente: Bases Técnicas Transversales.}{tab:rt-26}%
  {L{1.8cm} Y L{2cm} L{1.4cm} L{2.2cm} L{1.5cm} L{2cm}}%
  {\cab{Código} & \cab{Exigencia} & \cab{Carácter} & \cab{Cumple} & \cab{Componente} & \cab{Sección} & \cab{Evidencia}}
  RT-26.01 & Cada innovación se ubicará explícitamente en la arquitectura: qué capa la contiene, qué componentes la implementan y qué interfaces consume o expone. & Obligatorio &  &  &  &  \\
  RT-26.02 & Cada innovación identificará los paquetes de la estructura de descomposición del trabajo que la ejecutan y el mes del cronograma en que se materializa. & Obligatorio &  &  &  &  \\
  RT-26.03 & Las innovaciones de base tecnológica declararán el nivel de madurez de la tecnología con la escala utilizada y citarán las fuentes en norma APA 7.ª edición. & Obligatorio &  &  &  &  \\
  RT-26.04 & Cada innovación declarará su riesgo de adopción, su probabilidad, su impacto, la estrategia de mitigación y el plan de contingencia si no rinde lo esperado. & Obligatorio &  &  &  &  \\
  RT-26.05 & Cada innovación declarará su indicador de verificación con línea base, meta y momento de medición, y su impacto en inversión, costo operacional y beneficio esperado. & Obligatorio &  &  &  &  \\
  RT-26.06 & Las innovaciones que incorporen inteligencia artificial cumplirán íntegramente el Capítulo 18 de este documento. & Obligatorio &  &  &  &  \\
  RT-26.07 & Las innovaciones que modifiquen la arquitectura de seguridad requerirán su propio modelado de amenazas. & Obligatorio &  &  &  &  \\
  RT-26.08 & Se valorará que al menos una innovación sea verificable durante la marcha blanca de la Etapa 1, es decir, que su beneficio pueda medirse antes del mes 16. & Deseable &  &  &  &  \\
\end{tablalafrox}
